\chapter{الأوليات في الفترات القصيرة}
\label{ch:primes-short-intervals}

\section*{نطاق الفصل وحالته}

يدرس هذا الفصل مقدار ما يبقى من مبرهنة الأعداد الأولية عندما نستبدل الفترة الطويلة
\((0,x]\) بالفترة المتحركة
\[
(x,x+h],\qquad h=o(x).
\]
يثبت الفصل داخليًا لغة الفروق، ومبدأ نقل حد الخطأ، وضبط القوى الأولية العليا، والانتقال من \(\psi\) إلى \(\vartheta\) ثم \(\pi\)، وربط وجود أولي بحدود الفجوات. أما السجلات العددية العميقة فتُنقل من مصادرها الأصلية: صيغة Guth--Maynard التقاربية عند \(17/30+\varepsilon\)، وحد Baker--Harman--Pintz السفلي عند \(0.525\).

\noindent\textbf{حالة الفصل:} \textenglish{REVIEWED-CANDIDATE / NON-CITABLE}. اجتاز المتن بناء الدفعة والتدقيق الرياضي والمرجعي والمراجعة المستقلة، وأُغلقت ملاحظتاه الطباعيتان غير الحاجزتين. يبقى اعتماد المالك والدمج غير صادرين.

\subsection*{نواتج التعلم}

بنهاية الفصل يستطيع القارئ أن:
\begin{enumerate}
  \item يميز بين صيغة تقاربية، وحد سفلي من الرتبة الصحيحة، ومجرد الوجود، ونتيجة صحيحة لتقريبًا كل فترة.
  \item يكتب فروق \(\psi\)، و\(\vartheta\)، و\(\pi\) ويفهم ما تعده كل دالة.
  \item ينقل حد خطأ عالميًا معطى إلى فرق قصير من دون ادعاء نطاق لم يوفره المدخل.
  \item يزيل أثر القوى الأولية العليا عندما يكون \(h\) أكبر من \(x^{1/2+\eta}\).
  \item ينقل الصيغة التقاربية من \(\vartheta\) إلى \(\pi\) بحصار مباشر.
  \item يضع نتائج Hoheisel وIngham وHuxley وGuth--Maynard في ترتيبها التاريخي الصحيح.
  \item يشرح لماذا لا تُرتب نتيجتا Guth--Maynard وBaker--Harman--Pintz بمقارنة الأسين وحدهما.
\end{enumerate}

\section{أربعة أهداف مختلفة}

نستعمل، لأي دالة \(F\) على الأعداد الحقيقية الموجبة، الترميز
\[
\Delta_hF(x)=F(x+h)-F(x).
\]

\begin{definition}[أنظمة الفترات القصيرة]
\label{def:short-interval-regimes}
\resultid{ANT-DEF-19-01}
\provedhere
لتكن \(h=h(x)>0\) و\(h=o(x)\). نميز بين العبارات الآتية:
\begin{enumerate}
  \item \emph{صيغة تقاربية لكل \(x\)}:
  \[
  \Delta_h\pi(x)\sim\frac{h}{\log x}.
  \]
  \item \emph{حد سفلي من الرتبة الصحيحة لكل \(x\)}: وجود ثابت \(c>0\) بحيث
  \[
  \Delta_h\pi(x)\ge c\frac{h}{\log x}
  \]
  لكل \(x\) كبير.
  \item \emph{مجرد الوجود لكل \(x\)}:
  \[
  \Delta_h\pi(x)\ge1.
  \]
  \item \emph{نتيجة لتقريبًا كل \(x\)}: يسمح بمجموعة استثنائية كثافتها أو قياسها النسبي يؤول إلى صفر في المعنى المحدد في المبرهنة المستعملة.
\end{enumerate}
\end{definition}

قوة العبارات الثلاث الأولى تتناقص بهذا الترتيب. أما العبارة الرابعة فليست درجة رابعة على السلم نفسه: تغيير المكمم من ``كل \(x\)'' إلى ``تقريبًا كل \(x\)'' يمنع استعمالها لإثبات حكم نقطي عند كل موضع.

\begin{remark}
لا تكفي مبرهنة الأعداد الأولية النوعية
\(
\pi(x)\sim x/\log x
\)
لاستنتاج صيغة قصيرة عندما \(h=o(x)\). فقد يكون خطأ تقريب كل من \(\pi(x+h)\) و\(\pi(x)\) أكبر من الحد الرئيسي المتوقع \(h/\log x\). الفترات القصيرة مسألة عن دقة الخطأ، لا عن الحد الرئيسي العالمي وحده.
\end{remark}

\section{فروق دوال تشيبيشيف}

نذكر التعريفات
\[
\psi(x)=\sum_{n\le x}\Lambda(n),
\qquad
\vartheta(x)=\sum_{p\le x}\log p,
\qquad
\pi(x)=\sum_{p\le x}1.
\]

\begin{proposition}[هوية فرق فون مانغولت]
\label{prop:short-interval-psi-identity}
\resultid{ANT-ID-19-01}
\provedhere
لكل \(x\ge1\) و\(h>0\):
\[
\Delta_h\psi(x)
=
\sum_{x<n\le x+h}\Lambda(n).
\]
\end{proposition}

\begin{proof}
نطرح المجموعين المنتهيين
\[
\psi(x+h)=\sum_{n\le x+h}\Lambda(n),
\qquad
\psi(x)=\sum_{n\le x}\Lambda(n).
\]
تلغى الحدود ذات \(n\le x\)، وتبقى بالضبط الحدود التي تحقق \(x<n\le x+h\).
\end{proof}

إيجابية \(\Delta_h\psi(x)\) لا تثبت وحدها وجود أولي، لأن \(\Lambda(n)\) ترى القوى الأولية \(p^m\) أيضًا. هذه هي الفجوة التي يعالجها ضبط القوى العليا أدناه.

\section{مبدأ نقل حد الخطأ}

\begin{proposition}[نقل حد خطأ معطى إلى فرق قصير]
\label{prop:error-transfer-short-interval}
\resultid{ANT-PROP-19-01}
\provedhere
افترض أن
\[
\psi(t)=t+O(E(t))
\]
عند النقطتين \(t=x\) و\(t=x+h\). عندئذ
\[
\Delta_h\psi(x)
=
h+O\bigl(E(x+h)+E(x)\bigr).
\]
وبخاصة، إذا
\[
E(x+h)+E(x)=o(h),
\]
فإن
\[
\Delta_h\psi(x)\sim h.
\]
\end{proposition}

\begin{proof}
اكتب
\[
\psi(x+h)=x+h+O(E(x+h))
\]
و
\[
\psi(x)=x+O(E(x)).
\]
بطرح المعادلتين نحصل على الدعوى. والنتيجة التقاربية هي إعادة كتابة شرط أن يكون مجموع الخطأين \(o(h)\).
\end{proof}

\begin{remark}[حدود المبدأ]
هذه القضية لا تنتج أسًا عدديًا بذاتها. إذا كان المتاح فقط \(\psi(t)=t+o(t)\)، فلا يلزم أن يكون الفرق بين خطأي النقطتين \(o(h)\) عندما \(h=o(x)\). يجب أن يأتي النطاق القصير من حد خطأ كمي أو من مبرهنة قصيرة مقتبسة صراحة.
\end{remark}

\section{إزالة القوى الأولية العليا}

ضع
\[
R(t)=\psi(t)-\vartheta(t).
\]
من هوية القوى الأولية في الفصل التاسع:
\[
R(t)
=
\sum_{\substack{p^m\le t\\m\ge2}}\log p.
\]
إذن \(R\) غير سالبة وغير متناقصة؛ فكلما ازداد \(t\) لا تدخل إلا حدود غير سالبة جديدة.

\begin{lemma}[ضبط القوى العليا في فترة قصيرة]
\label{lem:higher-prime-powers-short-interval}
\resultid{ANT-LEM-19-01}
\provedhere
لكل \(x\ge2\) و\(h>0\):
\[
0\le
\Delta_h\psi(x)-\Delta_h\vartheta(x)
\le
R(x+h)
\ll
\sqrt{x+h}\log(x+h).
\]
وإذا كان \(h=o(x)\) ووجد \(\eta>0\) بحيث
\[
h\ge x^{1/2+\eta},
\]
فإن
\[
\Delta_h\psi(x)-\Delta_h\vartheta(x)=o(h).
\]
ومن ثم
\[
\Delta_h\psi(x)
=
h+o(h)
\quad\Longrightarrow\quad
\Delta_h\vartheta(x)=h+o(h).
\]
\end{lemma}

\begin{proof}
لأن \(R\) غير متناقصة وغير سالبة:
\[
0\le R(x+h)-R(x)\le R(x+h).
\]
لكن
\[
R(x+h)-R(x)
=
\Delta_h\psi(x)-\Delta_h\vartheta(x).
\]
وتعطي \textenglish{\texttt{ANT-LEM-09-02}} الحد
\[
R(t)\ll\sqrt t\log t.
\]
هذا يثبت الحصار الأول. وإذا كان \(h=o(x)\)، فإن \(x+h\sim x\)، ولذلك
\[
\frac{\sqrt{x+h}\log(x+h)}{h}
\ll
\frac{\sqrt{x}\log x}{x^{1/2+\eta}}
=
\frac{\log x}{x^\eta}
\longrightarrow0.
\]
فتكون مساهمة القوى العليا \(o(h)\)، وينتج الاستنتاج الأخير بالطرح.
\end{proof}

\section{الانتقال من \(\vartheta\) إلى \(\pi\)}

\begin{proposition}[حصار العد غير الموزون]
\label{prop:theta-to-pi-short-interval}
\resultid{ANT-PROP-19-02}
\provedhere
إذا كان \(x>1\) و\(h>0\)، فإن
\[
\frac{\Delta_h\vartheta(x)}{\log(x+h)}
\le
\Delta_h\pi(x)
\le
\frac{\Delta_h\vartheta(x)}{\log x}.
\]
وإذا كان \(h=o(x)\) و
\[
\Delta_h\vartheta(x)=h+o(h),
\]
فإن
\[
\Delta_h\pi(x)
\sim
\frac{h}{\log x}.
\]
\end{proposition}

\begin{proof}
لكل أولي \(p\) في \((x,x+h]\) لدينا
\[
\log x<\log p\le\log(x+h).
\]
بجمع هذه المتراجحات على الأوليات في الفترة نحصل على
\[
\log x\,\Delta_h\pi(x)
\le
\Delta_h\vartheta(x)
\le
\log(x+h)\,\Delta_h\pi(x),
\]
ثم نعيد ترتيبها. وإذا كان \(h=o(x)\)، فإن
\[
\log(x+h)=\log x+\log\!\left(1+\frac hx\right)
=\log x+o(1),
\]
وبالتالي \(\log(x+h)/\log x\to1\). نعوض الصيغة المفترضة في الحصار ونطبق مبرهنة الحصر.
\end{proof}

\begin{remark}
هذا الانتقال محلي ولا يحتاج إلى حد خطأ عالمي جديد أو إلى استعمال دوري للجمع الجزئي. لكنه يحتاج إلى صيغة لـ\(\vartheta\) ذات رتبة خطأ أصغر من \(h\).
\end{remark}

\section{المسار التاريخي: من Hoheisel إلى Huxley}

أثبت Hoheisel لأول مرة أن مبرهنة الأعداد الأولية تبقى صحيحة في فترات طولها \(x^{1-\delta}\) لثابت ما \(\delta>0\) \cite{Hoheisel1930}. لا نسجل هنا قيمة عددية لـ\(\delta\)، لأن النص الأصلي لم يكن متاحًا رقميًا في تدقيق المشروع وتعارضت بعض المصادر الثانوية في نقلها.

ربط Ingham النطاق القصير بحدود دالة زيتا على الخط الحرج. في صياغته، إذا
\[
\zeta\!\left(\frac12+it\right)=O(t^c),
\]
ظهر الأس
\[
\frac{1+4c}{2+4c},
\]
وأعطت المدخلات الكلاسيكية المرحلة \(5/8\) \cite[Theorem~1]{Ingham1937}. ثم خفض Huxley السجل التقاربي التاريخي إلى \(7/12\) \cite{Huxley1972}. هذه النتائج مقتبسة ومشروحة؛ لا يعيد الفصل بناء تقديرات كثافة الأصفار التي تحملها.

\section{السجل التقاربي الحديث}

\begin{theorem}[Guth--Maynard]
\label{thm:guth-maynard-short-intervals}
\resultid{ANT-THM-19-01}
\citedresult[\cite{GuthMaynard2026}، Corollary~1.3]
لكل \(\varepsilon>0\)، وبانتظام في النطاق
\[
x^{17/30+\varepsilon}\le h\le x^{0.99},
\]
لدينا
\[
\pi(x+h)-\pi(x)
=
\frac{h}{\log x}
+O_\varepsilon\!\left(
h\exp\!\bigl(-\sqrt[4]{\log x}\bigr)
\right)
\]
عندما يكون \(x\) كبيرًا بما يكفي.
\end{theorem}

يعتمد البرهان الأصلي على تقديرات جديدة للقيم الكبيرة لمتعددات حدود ديريشليه، وتقديرات كثافة للأصفار، ومنطقة Vinogradov--Korobov الخالية من الأصفار، وصيغة صريحة مع موازنة دقيقة للمعلمات. يعرض الفصل هذه السلسلة بوصفها \textenglish{CITED-CORE / EXPLANATION ONLY}، ولا يعيد تسميتها \textenglish{PROVED-HERE}.

وبما أن
\[
\frac{17}{30}=0.566\ldots<\frac{7}{12}=0.583\ldots,
\]
فإن نتيجة Guth--Maynard حسنت السجل التقاربي السابق. لكن هذا لا يجعلها تلقائيًا ``أقوى من كل نتيجة ذات أس أصغر''؛ يجب أولًا مقارنة نوع النتيجة ومكمماتها.

\section{فترة أقصر بنتيجة أضعف نوعيًا}

\begin{theorem}[Baker--Harman--Pintz]
\label{thm:bhp-short-interval-lower-bound}
\resultid{ANT-THM-19-02}
\citedresult[\cite{BakerHarmanPintz2001}، خاتمة ص.~561]
لكل \(x\) كبير بما يكفي:
\[
\pi\!\left(x+x^{0.525}\right)-\pi(x)
>
\frac{9}{100}\frac{x^{0.525}}{\log x}.
\]
\end{theorem}

تضمن هذه المبرهنة وجود أولي في فترة أقصر من النطاق التقاربي السابق، وتعطي حدًا سفليًا من الرتبة الصحيحة. لكنها لا تثبت
\[
\pi(x+x^{0.525})-\pi(x)
\sim
\frac{x^{0.525}}{\log x}.
\]
إذن لا يجوز ترتيب نتيجتي Guth--Maynard وBaker--Harman--Pintz بمقارنة \(17/30\) و\(21/40\) وحدهما: الأولى أقوى نوعيًا في فترة أطول، والثانية أضعف نوعيًا في فترة أقصر.

\begin{remark}[ادعاء \(0.52\)]
تدعي المسودة \cite{LiShortIntervalsPreprint} وجود أوليات في فترات ذات أس \(0.52\). وفي تاريخ القطع الأدبي للمشروع كانت النسخة \textenglish{arXiv:2308.04458v8} غير مثبتة كسجل محكّم؛ لذلك تبقى \textenglish{QUARANTINED-PREPRINT}، ولا تحمل معرّف نتيجة في الموسوعة، ولا تستبدل السجل المعتمد.
\end{remark}

\section{من وجود أولي إلى فجوة بين أوليين متتاليين}

\begin{corollary}[حد الفترة يعطي حد الفجوة]
\label{cor:short-interval-to-prime-gap}
\resultid{ANT-COR-19-01}
\provedhere
افترض وجود \(0<\theta<1\) و\(x_0\) بحيث تحوي الفترة
\[
(x,x+x^\theta]
\]
أوليًا لكل \(x\ge x_0\). عندئذ، لكل \(n\) كبير بما يكفي:
\[
p_{n+1}-p_n\le p_n^\theta.
\]
\end{corollary}

\begin{proof}
خذ \(x=p_n\ge x_0\). بالفرض يوجد أولي \(q\) يحقق
\[
p_n<q\le p_n+p_n^\theta.
\]
وبتعريف \(p_{n+1}\) بوصفه أصغر أولي أكبر من \(p_n\)، لدينا \(p_{n+1}\le q\). لذلك
\[
p_{n+1}-p_n
\le q-p_n
\le p_n^\theta.
\]
\end{proof}

\begin{remark}[اتجاه الاستنتاج]
العكس ليس آليًا بالصيغة نفسها: حد مكتوب بدلالة \(p_n\) يحتاج إلى اختيار الأولي السابق لـ\(x\) وضبط المقارنة بين \(x\) و\(p_n\). لذلك نسجل هنا الاتجاه المستعمل فقط.
\end{remark}

\section{ما الذي تصنعه الصيغة الصريحة؟}

في الصورة الإرشادية، تربط الصيغة الصريحة فرق \(\psi\) بمجموع على أصفار زيتا:
\[
\Delta_h\psi(x)
=
h-
\sum_\rho
\frac{(x+h)^\rho-x^\rho}{\rho}
+\text{حدود القطع والخطأ}.
\]
لا تكفي كتابة هذه الصيغة رسميًا. يجب اختيار ارتفاع القطع، وضبط عدد الأصفار قرب خطوط رأسية مختلفة، واستعمال منطقة خالية من الأصفار، ثم موازنة الخطأ مع \(h\). تاريخيًا حسنت تقديرات كثافة الأصفار هذا التوازن، وفي عمل Guth--Maynard جاءت قفزة \(17/30\) من تقديرات قيم كبيرة جديدة تدخل في آلة الكثافة.

هذا الشرح لا يسد دين بيرون وتحويل المسار في الفصول السابقة، ولا يدعي إعادة برهان النتيجة الحديثة. الغرض منه تحديد موضع العمق: الهويات المحلية سهلة نسبيًا، أما الوصول إلى أس قصير موحد لكل \(x\) فهو محمول على معلومات دقيقة عن أصفار زيتا ومتعددات حدود ديريشليه.

\section{تقريبًا كل فترة}

قد تسمح النتائج المتوسطية بفترات أقصر إذا استثنينا مجموعة صغيرة من مواضع البداية. لكن لا يجوز إدخال هذه النتائج في برهان حكم لكل \(x\). لهذا يفصل الفصل بين المكممين:
\[
\text{لكل \(x\)}
\qquad\text{و}\qquad
\text{لتقريبًا كل \(x\)}.
\]
أي نتيجة متوسطية تضاف لاحقًا يجب أن تصرح بدقة بمجال المتوسط، وقياس المجموعة الاستثنائية، وانتظام الثوابت، وألا تُستعمل لإثبات \textenglish{ANT-THM-19-01} أو \textenglish{ANT-THM-19-02}.

\section{أخطاء شائعة}

\begin{enumerate}
  \item طرح صيغتين من نوع \(x+o(x)\) ثم افتراض أن الفرق يساوي \(h+o(h)\) بلا حد خطأ كمي.
  \item استنتاج وجود أولي من \(\Delta_h\psi>0\) من دون إزالة القوى الأولية العليا.
  \item وصف حد سفلي من الرتبة الصحيحة بأنه صيغة تقاربية.
  \item تحويل نتيجة ``تقريبًا كل \(x\)'' إلى نتيجة ``كل \(x\)''.
  \item وصف \(7/12\) بأنه السجل التقاربي الحالي بعد نتيجة Guth--Maynard.
  \item اعتماد مسودة \(0.52\) كما لو كانت ورقة محكّمة.
  \item مقارنة نتيجتين بمقارنة الأسين فقط مع تجاهل نوع الحكم.
\end{enumerate}

\section{تمارين}

\subsection{المستوى الأول}
\begin{enumerate}
  \item أثبت هوية \(\Delta_h\vartheta(x)=\sum_{x<p\le x+h}\log p\).
  \item تحقق من أن \(R(t)=\psi(t)-\vartheta(t)\) غير متناقصة.
  \item أثبت أن \(\log(x+h)\sim\log x\) إذا \(h=o(x)\).
  \item فسر لماذا تستلزم الصيغة التقاربية حدًا سفليًا من الرتبة الصحيحة.
\end{enumerate}

\subsection{المستوى الثاني}
\begin{enumerate}
  \item إذا كان \(E(t)=t\exp(-c\sqrt{\log t})\)، استخرج شرطًا كافيًا على \(h\) كي يعطي مبدأ نقل الخطأ \(\Delta_h\psi(x)\sim h\).
  \item أثبت أن \(\sqrt{x}\log x=o(x^{1/2+\eta})\) لكل \(\eta>0\).
  \item أعد برهان الانتقال \(\vartheta\to\pi\) باستعمال جمع جزئي، ثم قارن طول البرهان بالحصار المباشر.
  \item بين أن حدًا سفليًا موجبًا لـ\(\Delta_h\vartheta\) يثبت وجود أولي مباشرة.
\end{enumerate}

\subsection{المستوى الثالث}
\begin{enumerate}
  \item حدد بدقة أين تفشل محاولة استنتاج نتيجة قصيرة من PNT النوعية وحدها.
  \item اشرح لماذا يحتاج عكس \textenglish{\texttt{ANT-COR-19-01}} إلى ضبط إضافي.
  \item ارسم خريطة اعتماد تبدأ من صيغة صريحة مقطوعة وتنتهي بحد على \(\Delta_h\psi\)، وحدد المدخلات التي تبقى مقتبسة.
  \item قارن منطقيًا بين مبرهنة Guth--Maynard ومبرهنة Baker--Harman--Pintz من حيث طول الفترة، ونوع النتيجة، والمكمم.
\end{enumerate}

\section*{حالة نتائج المسودة}

\begin{center}
\begin{tabular}{ll}
\texttt{ANT-DEF-19-01} & \texttt{AUTHORED-DRAFT / DEFINITION}\\
\texttt{ANT-ID-19-01} & \texttt{AUTHORED-DRAFT / PROVED-HERE}\\
\texttt{ANT-PROP-19-01} & \texttt{AUTHORED-DRAFT / PROVED-HERE}\\
\texttt{ANT-LEM-19-01} & \texttt{AUTHORED-DRAFT / PROVED-HERE}\\
\texttt{ANT-PROP-19-02} & \texttt{AUTHORED-DRAFT / PROVED-HERE}\\
\texttt{ANT-THM-19-01} & \texttt{AUTHORED-DRAFT / CITED-EXPLAINED}\\
\texttt{ANT-THM-19-02} & \texttt{AUTHORED-DRAFT / CITED}\\
\texttt{ANT-COR-19-01} & \texttt{AUTHORED-DRAFT / PROVED-HERE}
\end{tabular}
\end{center}

جميع النتائج الثماني مؤلفة، لكنها تبقى \textenglish{NON-CITABLE}. لا تصبح \textenglish{ACTIVE} إلا بعد نجاح البناء، والتدقيق الرياضي والمرجعي، والمراجعة المستقلة للمتن، وقرار المالك.

\section{خلاصة الفصل}

الفترات القصيرة تكشف الفرق بين معرفة الكثافة العالمية للأوليات ومعرفة انتظامها المحلي. الهويات التي تنقل من \(\psi\) إلى \(\vartheta\) ثم \(\pi\) قصيرة وواضحة، لكن المدخل الذي يجعل الخطأ أصغر من طول الفترة هو موضع العمق الحقيقي. تاريخيًا حملت كثافة الأصفار هذا العبء من Hoheisel إلى Huxley، ثم نقل Guth--Maynard السجل التقاربي إلى \(17/30+\varepsilon\). وفي اتجاه مختلف، أعطى Baker--Harman--Pintz وجودًا وحدًا سفليًا من الرتبة الصحيحة في فترة أقصر عند \(0.525\). الفصل الصحيح لا يخلط هذين النوعين، ولا يحول مسودة أحدث إلى سجل معتمد.
